% Page 010 — page imprimée 6
% Suite de l'interrogatoire de Jacques Crégut
% Reconstruction éditable en LaTeX.


Desmaze faisant recrue dans le temps qu’il travaillait étant derrière le métier de tisserand à la maison du sieur Laval, le fils du Sr. Desmaze accompagné de deux de ses soldats et la recrue seraient venus le prendre en lui disant qu’il était enrôlé avec lui. Etant resté deux ou trois jours à la maison du Sr. Desmaze, voyant qu’il était enrôlé par force, il aurait déserté pour aller à Anduze travailler de son métier de tisserand, au commencement du mois de Février dernier.

Pour aller à \textbf{Anduze} il serait passé par \textbf{Mandagout, Taleyrac, Bonnerrier, les Plantier de Saint Michel de Fonfouillouse}, par \textbf{Saumane} où il aurait rencontré la troupe de fanatiques dont le chef était le nommé Castanet de \textbf{Massevaques}.

C’était un petit homme assez gros, portant une perruque blonde, habillé de minime \textit{(tunique brune)} doublée de rouge, portant sur le chapeau un crêpe, disant qu’il faisait le deuil de son père.

Son lieutenant était le nommé Laroze, tisserand de son métier, ne sachant pas son nom ni sa demeure. Etant d’une grande taille, portant une perruque noire, maigre de visage, la troupe de fanatiques était composée d’environ 140 hommes et en ayant 40 d’armés. Le Sr. Castanet menait une jeune fille faisant la prophétesse, âgée d’environ 15-16 ans qu’on dit être de Saint \textbf{André de Valborgne}.

Le Sr. Castanet a obligé Crégut à le suivre, lui ayant baillé un fusil, une chemise blanche qui est la même qu’il porte présentement, et un galon d’or pour border son chapeau comme il fit.

L’ayant logé le même jour à \textbf{Saumane} chez le nommé Bastide avec trois autres inconnus. Et aussi le reste de la troupe, quatre par quatre, chez les habitants de Saumane, le Sr. Crégut n’ayant connu aucun de la troupe.

Thomas Daumessenne, qui avait déserté de la compagnie du Sr. Deleveret, ayant soupé à Saumane ; ils en seraient partis ayant marché toute la nuit ne sachant où ils furent... Le lendemain au château de l’Oum ( ?) le nommé Etienne fils du nommé Fabre dit Causses Vertes de \textbf{Pracoustal}, paroisse d’\textbf{Aulas} fut les joindre ayant resté dans la troupe du Sr. Castanet depuis ; ce fait environ 15 jours qu’il les quitta. Le Sr. Castanet les faisaient marcher de nuit le plus souvent et le jour, dans le pays désert, les faisaient reposer.

Pellenc ( ?) qui leur portait les vivres n’en connaissant aucun, obéissant au commandement de Castanet, ne restant presque jamais au même endroit et ne sachant jamais où Castanet les conduirait. Le Sr. Castanet ayant beaucoup d’argent qu’il faisait rançonner dans plusieurs endroits.

Castanet ayant prêché deux fois dans l’église de \textbf{St. Michel de Fonfouillouse} suivi de sa troupe. Beaucoup des habitants de la paroisse et des environs y ayant assisté ; n’en connaissait aucun.

Le Sr. Castanet ayant fait brusler trois maisons à \textbf{St. André de Valborgne}, ne sachant à qui elles appartenaient. Il fit ensuite le soir, dans la nuit, une assemblée en dessous de St. André, dans un pré à deux portées de fusil, étant venues plusieurs personnes de St. André et des environs qu’il ne connut point.

Il fut au Pompidou et se plaignait que les habitants de \textbf{Fraissinet de Fourques} l’avaient menacé de brûler sa maison.

Il fut à \textbf{Fraissinet} avec sa troupe qui a grossi jusqu’à 200. Les 2/3 armés ou en ayant jamais plus, en commandant de brûler et tuer, disant que ceux qui ne le feraient pas, il leur tirerait dessus. Les habitants de Fraissinet les ayant vu venir se seraient réfugiés à la maison où ils

